\documentclass[11pt]{article}
\usepackage{amsmath,amssymb,amsthm}
\usepackage[margin=1.2in]{geometry}
\usepackage{hyperref}

\newcommand{\Tp}{T_p}
\newcommand{\Ta}{T_a}
\newcommand{\Te}{T_e}
\newcommand{\kk}{k}
\newcommand{\tauu}{\tau}
\renewcommand{\det}{\mathrm{det}}
\newcommand{\Cov}{\mathrm{Cov}}

\title{Loss Functions for Active and Passive Diffusion:\\
Comparison with CLD and the $\tauu\to 0$ Limit}
\author{}
\date{}

\begin{document}
\maketitle

\section{Setup}

We consider two coupled stochastic variables: image pixels $x$ and active noise $\eta$,
evolving under the forward SDEs
\begin{align}
dx &= (-\kk x + \eta)\,dt + \sqrt{2\Tp}\,dW_x, \label{eq:sde_x}\\
d\eta &= -\frac{\eta}{\tauu}\,dt + \sqrt{\frac{2\Ta}{\tauu^2}}\,dW_\eta, \label{eq:sde_eta}
\end{align}
with $\Tp = 10^{-3}$, $\Ta = 1$, and persistence time $\tauu$.
The passive baseline is the single-variable OU process
\begin{equation}
dx = -\kk x\,dt + \sqrt{2\Te}\,dW_x.
\end{equation}

\section{Marginal Distribution and the $\tauu\to 0$ Limit}

The joint process $(x_t, \eta_t)$ given $(x_0, \eta_0)$ is Gaussian with means
\begin{equation}
\mu_x = e^{-\kk t} x_0 + b\,\eta_0, \qquad
\mu_\eta = e^{-t/\tauu}\,\eta_0,
\end{equation}
where $b = (e^{-t/\tauu} - e^{-\kk t})/(\kk - 1/\tauu) \to 0$ as $\tauu\to 0$,
and covariance matrix $\Sigma = \bigl[\begin{smallmatrix} M_{11} & M_{12} \\ M_{12} & M_{22} \end{smallmatrix}\bigr]$ with
\begin{equation}
M_{11} \xrightarrow{\tauu\to 0} \frac{\Tp + \Ta}{\kk}(1 - e^{-2\kk t}), \qquad
M_{12} \xrightarrow{\tauu\to 0} \Ta, \qquad
M_{22} = \frac{\Ta}{\tauu} \to \infty.
\end{equation}
The marginal variance of $x_t$ therefore tends to $(\Tp+\Ta)/\kk\,(1-e^{-2\kk t})$,
which matches the passive model with $\Te = \Tp + \Ta$.
In particular, when $\Tp\to 0$ and $\Ta = \Te$, the marginal distribution of $x_t$ is
identical to the passive model with temperature $\Te$.

\section{Score-Matching Loss Functions}

\subsection{Passive loss}

The passive model uses true score $\nabla_x \log p_t(x_t) = -(x_t - \mu_x)/M_{11}$
where $M_{11} = (T_e/k)(1 - e^{-2kt})$.
The loss is implemented via \texttt{mse\_loss} with two variants depending on GPU mode:
\begin{align}
\mathcal{L}_{\text{passive}}^{\text{AMP}} &= M_{11} \cdot \bigl(F_x - \nabla_x \log p_t\bigr)^2,
\label{eq:passive_amp}\\
\mathcal{L}_{\text{passive}}^{\text{non-AMP}} &= M_{11}^2 \cdot \bigl(F_x - \nabla_x \log p_t\bigr)^2.
\label{eq:passive_noamp}
\end{align}
These arise from \texttt{mse\_loss($\sigma F_x$,\; $\sigma \nabla_x \log p_t$)}
and \texttt{mse\_loss($\sigma^2 F_x$,\; $\sigma^2 \nabla_x \log p_t$)} respectively,
with $\sigma^2 = M_{11}$.

\paragraph{CLD passive (VPSDE).}
CLD's passive baseline uses variance-preserving noise schedule $\beta(t)$,
with $M_{11} = 1 - e^{-\int_0^t \beta(s)\,ds}$ and loss
\begin{equation}
\mathcal{L}_{\text{CLD-passive}} = \tfrac{1}{2}\beta(t) \cdot M_{11}
\cdot \bigl(F_x - \nabla_x \log p_t\bigr)^2.
\label{eq:cld_passive}
\end{equation}
This follows from \texttt{(score / noise\_multiplier $-\,\varepsilon)^2 \times$ multiplier}
with noise\_multiplier $= -1/\sigma$ and multiplier $= \beta(t)/2$.

\paragraph{Comparison.}
Equations~\eqref{eq:passive_amp} and~\eqref{eq:cld_passive} have the same $M_{11}$
weighting; the only difference is the extra $\tfrac{1}{2}\beta(t)$ factor in CLD, which
reflects its time-varying noise schedule (in the user's model $\beta$ is constant, so
this is just a fixed prefactor absorbed into the learning rate).
Equation~\eqref{eq:passive_noamp} carries an additional $M_{11}$ factor; since
$M_{11} \in (0, T_e/k]$ is bounded and $\mathcal{O}(1)$, this changes the relative
emphasis across diffusion times but not the minimizer.

\subsection{Active loss (joint score matching)}
The active model learns the full joint score
$\nabla_{(x,\eta)}\log p_t(x_t,\eta_t) = -\Sigma^{-1}(x_t - \mu_x, \eta_t - \mu_\eta)^\top$:
\begin{equation}
\mathcal{L}_{\text{active}} =
\Tp \cdot \det\Sigma \cdot \bigl(F_x - \nabla_x \log p_t\bigr)^2
+\Ta \cdot \det\Sigma \cdot \bigl(F_\eta - \nabla_\eta \log p_t\bigr)^2,
\label{eq:active_loss}
\end{equation}
where
\begin{align}
\nabla_x \log p_t &= \frac{-M_{22}(x_t - \mu_x) + M_{12}(\eta_t - \mu_\eta)}{\det\Sigma},
\label{eq:x_score}\\
\nabla_\eta \log p_t &= \frac{-M_{11}(\eta_t - \mu_\eta) + M_{12}(x_t - \mu_x)}{\det\Sigma}.
\label{eq:eta_score}
\end{align}
In practice, the loss is implemented by introducing
$\mathbf{F}_x = \sqrt{\det\Sigma}\,F_x$ and $\mathbf{F}_\eta = \sqrt{\det\Sigma}\,F_\eta$,
so that the targets become $\mathcal{O}(1)$ quantities independent of $\tauu$.

\subsection{CLD loss (conditional score matching)}

Critically-Damped Langevin Diffusion \cite{dockhorn2022cld} augments data $x$ with a
velocity variable $v$ and trains only the conditional score $\nabla_v \log p_t(v_t\mid x_t)$:
\begin{equation}
\mathcal{L}_{\text{CLD}} = \beta(t)\,f \cdot \frac{\det\Sigma}{M_{11}} \cdot
\bigl(F_v - \nabla_v \log p_t\bigr)^2,
\end{equation}
where $\det\Sigma/M_{11} = \mathrm{Var}(v_t\mid x_t)$ is the conditional variance of $v$
given $x$.  Since $\nabla_v \log p(x,v) = \nabla_v \log p(v\mid x)$ (as $p(x)$ does not
depend on $v$), the joint and conditional $v$-scores coincide, so CLD needs no separate
$x$-score term.  The weighting $\det/M_{11}$ (conditional variance) is more natural than
the joint determinant $\det$ used in~\eqref{eq:active_loss}; the ratio is the
bounded, $\mathcal{O}(1)$ quantity $M_{11}$.

\section{Behavior of the Active Loss as $\tauu\to 0$}

As $\tauu\to 0$ with $\Tp = 10^{-3}$, $\Ta = 1$ fixed, we have
$\det\Sigma \approx M_{11}\cdot\Ta/\tauu$.

\paragraph{$x$-score term.}
\begin{equation}
\Tp\cdot\det\Sigma \;\approx\; \frac{\Tp\,\Ta\,M_{11}}{\tauu}
\;=\; \frac{\Tp}{\tauu}\times \Ta\,M_{11}.
\end{equation}
Since $\nabla_x\log p_t\to$ (passive score with $\Te=\Tp+\Ta$), this term becomes
$(\Tp/\tauu)$ times the passive loss---an infinitely amplified version of the passive
score-matching objective.  The model $F_x$ is therefore driven strongly to the passive
score for any $\Tp > 0$, no matter how small.

\paragraph{$\eta$-score term and $x$-information.}
The target for $F_\eta$ contains an $x$-dependent coupling:
\begin{equation}
\nabla_\eta \log p_t \;\approx\;
-\frac{({\eta_t - \mu_\eta})}{\Ta/\tauu}
+\underbrace{\frac{\tauu\,(x_t - \mu_x)}{M_{11}}}_{\to\,0}.
\label{eq:eta_score_approx}
\end{equation}
The coupling term vanishes as $\tauu\to 0$, yet the loss gradient for it is
\begin{equation}
\Ta\cdot\det\Sigma\cdot\frac{\tauu}{M_{11}}
\;\approx\;
\Ta\cdot M_{11}\cdot\frac{\Ta}{\tauu}\cdot\frac{\tauu}{M_{11}} = \Ta^2 = \mathcal{O}(1).
\end{equation}
The det$\,\Sigma\to\infty$ weighting exactly compensates the $\tauu\to 0$ shrinkage of the
coupling, keeping the gradient signal finite and well-conditioned at all $\tauu$.

\section{Adiabatic Elimination: Active Score Reduces to Passive}

The reverse-time SDE for $\eta$ (running backward from $T$ to $0$) is
\begin{equation}
d\eta_{\rm rev} = \left[\frac{\eta}{\tauu}
+ \frac{2\Ta}{\tauu^2}\nabla_\eta\log p_t\right]ds + \sqrt{\frac{2\Ta}{\tauu^2}}\,d\bar W.
\label{eq:rev_sde_eta}
\end{equation}
Setting the mean drift to zero defines the quasi-stationary point $\eta^*$:
\begin{equation}
\frac{\eta^*}{\tauu} + \frac{2\Ta}{\tauu^2}\,\nabla_\eta\log p_t(x,\eta^*) = 0
\quad\Longrightarrow\quad
\nabla_\eta\log p_t(x,\eta^*) = -\frac{\tauu\,\eta^*}{2\Ta}.
\label{eq:qs_cond}
\end{equation}
We now substitute the \emph{exact} score formula \eqref{eq:eta_score} into \eqref{eq:qs_cond},
with $\mu_\eta\approx 0$:
\begin{equation}
\frac{-M_{11}\,\eta^* + M_{12}(x-\mu_x)}{\det\Sigma}
= -\frac{\tauu\,\eta^*}{2\Ta}.
\end{equation}
Multiplying through by $\det\Sigma$ and using $\det\Sigma = M_{11}M_{22} - M_{12}^2$
with the \emph{exact} fluctuation--dissipation relation $M_{22} = \Ta/\tauu$:
\begin{equation}
\tauu\,\frac{\det\Sigma}{2\Ta}
= \tauu\,\frac{M_{11}M_{22}-M_{12}^2}{2\Ta}
= \frac{M_{11}}{2} - \frac{\tauu\,M_{12}^2}{2\Ta}.
\end{equation}
Collecting the $\eta^*$ terms gives
\begin{equation}
M_{12}(x-\mu_x)
= \eta^*\!\left(M_{11} - \frac{M_{11}}{2} + \frac{\tauu\,M_{12}^2}{2\Ta}\right)
= \eta^*\,\frac{T_a M_{11} + \tauu\,M_{12}^2}{2\Ta},
\end{equation}
and therefore
\begin{equation}
\boxed{\eta^* = \frac{2\Ta\,M_{12}\,(x-\mu_x)}{\Ta\,M_{11} + \tauu\,M_{12}^2}}.
\label{eq:eta_star}
\end{equation}

\paragraph{Closing the loop: substitute $\eta^*$ into the $x$-score \eqref{eq:x_score}.}
With $\mu_\eta\approx 0$:
\begin{align}
\nabla_x\log p_t\big|_{\eta=\eta^*}
&= \frac{-M_{22}(x-\mu_x) + M_{12}\,\eta^*}{\det\Sigma}
= \frac{(x-\mu_x)}{\det\Sigma}
\left[-M_{22} + \frac{2\Ta\,M_{12}^2}{\Ta\,M_{11}+\tauu\,M_{12}^2}\right].
\notag
\end{align}
Combining the bracket over a common denominator and using $M_{11}M_{22}=\det\Sigma+M_{12}^2$:
\begin{align}
-M_{22}(\Ta M_{11}+\tauu M_{12}^2)+2\Ta M_{12}^2
&= -\Ta M_{11}M_{22} - \tauu M_{22}M_{12}^2 + 2\Ta M_{12}^2
\notag\\
&= -\Ta(\det\Sigma+M_{12}^2) + M_{12}^2(\underbrace{2\Ta - \tauu M_{22}}_{=\,\Ta,\;\text{since } M_{22}=\Ta/\tauu})
\notag\\
&= -\Ta\,\det\Sigma.
\label{eq:cancel}
\end{align}
The cancellation $2\Ta - \tauu M_{22} = \Ta$ (using $M_{22}=\Ta/\tauu$ \emph{exactly}) is key.
Substituting back:
\begin{equation}
\nabla_x\log p_t\big|_{\eta=\eta^*}
= -\frac{\Ta\,(x-\mu_x)}{\Ta\,M_{11} + \tauu\,M_{12}^2}
\;\xrightarrow{\tauu\to 0}\; -\frac{x-\mu_x}{M_{11}}.
\label{eq:x_score_passive}
\end{equation}
Since $M_{11}\to(\Tp+\Ta)/\kk\,(1-e^{-2\kk t})=(\Te/\kk)(1-e^{-2\kk t})$ as $\tauu\to 0$,
this is exactly the passive score for temperature $\Te=\Tp+\Ta$, valid for any $\Tp\geq 0$
including $\Tp=0$.

\medskip
\noindent\textbf{Conclusion.}
Using the exact $\eta$-score \eqref{eq:eta_score} in the quasi-stationary condition
\eqref{eq:qs_cond} yields $\eta^*$~\eqref{eq:eta_star}, and substituting into
\eqref{eq:x_score} recovers the passive score \eqref{eq:x_score_passive} via the
exact cancellation \eqref{eq:cancel} (which follows from the fluctuation--dissipation
relation $M_{22}=\Ta/\tauu$).  The $\det\Sigma$ weighting in the active loss keeps
the gradient for the $x$-coupling in $F_\eta$ at $\mathcal{O}(\Ta^2)=\mathcal{O}(1)$
for all $\tauu$, ensuring the network learns $\nabla_\eta\log p_t$ accurately and
the passive limit is recovered.

\begin{thebibliography}{1}
\bibitem{dockhorn2022cld}
T.~Dockhorn, A.~Vahdat, and K.~Kreis,
\textit{Score-Based Generative Modeling with Critically-Damped Langevin Diffusion},
ICLR 2022.
\end{thebibliography}

\end{document}
